\scanpage{10}
\deleted{qui est aussi la relation (unique, \inserted{à unité près,} on la vérifie aussi bien à priori) vérifiée par} Puisque $1,y,x,y^2,xy,y^3$ forment une base de $\mathcal E_6$ (adaptée à sa filtration) on aura donc une équation
\[
\varphi(x,y)=\boxed{x^2+ay^3+bxy+cy^2+dx+ey+f=0}
\]

\[
\Delta(a,b,c,d,e,f)\in L^{\otimes12}
\]

qui est d'ailleurs une équation \unsure{valable|soluble} dans $\mathcal E_6$ considéré comme composante de $\Gamma_{*\mathrm{\unsure{rel}}}(f,\underline L^{\otimes})\simeq\coprod_{i\geq0}\mathcal E_{3i}$, satisfaite entre les sections $x,y$ de $\mathcal E_3$ \deleted{\unsure{et $x$ de $\mathcal E_3$|$x_1,x_2$} \unsure{$y$ de $\mathcal E_2$|$x_3$}}, \inserted{à valeurs en $\mathcal E_{\unsure{3.3|9}}$ \deleted{\unsure{resp.}}} \deleted{Cette équation} Si on désigne par $z$ la section $1$ de $\mathcal E_0,=\mathcal E_1\subset\mathcal E_3$, et qu'on \deleted{\unsure{la}} écrit \inserted{\unsure{une|ici} \deleted{la} équation} sous la forme
\[
\Phi(x,y,z)=x^2z+ay^3\deleted{z}+bxyz+cy^2z+dxz^2+fz^3=0
\]
% slanted
\margin{détailler, et éventuellement généraliser}
on \uline{constate alors} que $\Phi$ est la forme de degré 3 qui définit \deleted{\unsure{alors}} $C$ dans $P(\mathcal E_3)\simeq\mathbb P^3_S$. En effet si $\Phi'$ est la forme qui définit $C$, alors $\Phi$ est \uline{un multiple de $\Phi'$} et \deleted{\unsure{d'un coef.}} c'est un multiple par quelque chose qui est inversible, puisque $\Phi$ est ``non nulle'' en tt pt de $S$ --- un de ses coefficients c'est 1 ---.

% slanted
\margin{préciser ce pt cf intersections complètes}
% archive sheet number: 9
